\subsection{Confirmation with intermediate progress capture}
\label{sec:progress-confirmation}
A separate confirmation uses the revised \texttt{isolated-progress-v2}
protocol on the same 39 graphs, three parameter pairs, three repetitions,
and 30-second external budget. The complete file
\path{general_progress_v2_20261007T062406751526Z.csv} contains
\ProgressObservations{} observations. An earlier interrupted attempt is
excluded rather than merged into the completed run.
SAT transmits completed span decisions; Gurobi transmits validated
incumbent updates and conservative integer dual bounds. The parent checks
received intervals and retains the last accepted update at interruption.
Progress timestamps measure parent receipt, not native solver discovery.
Callback, transport, and intermediate validation overhead are included in
wall time; v1 and v2 timing observations are not pooled.

All saved labelings and progress ledgers pass the audit, with no
contradictory OPT values or intervals across backends and repetitions.
As before, witness validation does not independently verify UNSAT or dual
certificates. The data support the coverage results in
Table~\ref{tab:progress-coverage}. Both backends satisfy the all-three-OPT
criterion on \ProgressJointStable{} instances.

\begin{table}[htbp]
\centering\small
\caption{Separate v2 confirmation: optimality coverage and retained feasible
solutions. $^*$FEASIBLE observations whose upper bound improves on the
initial greedy span. Improvements do not establish optimality.}
\label{tab:progress-coverage}
\input{\ProgressReportDir/coverage.tex}
\end{table}

Pooling only the completed v2 run's solver bounds across repetitions and
backends closes \ProgressClosed{} instance intervals and leaves
\ProgressOpen{} open, all in the ER family. This post-hoc combination is
not attributed to a single timed solver run. The open instances give
\ProgressOpenPairs{} paired observations, not independent graph samples.
Table~\ref{tab:progress-bounds} distinguishes larger lower bounds from
smaller feasible upper bounds.

\begin{table}[htbp]
\centering\small
\caption{Paired retained bounds in v2. Each pair matches graph, $(h,k)$,
and repetition; ``open'' is defined by pooling this run's intervals.
Counts for all instances and the open subset must not be added together.}
\label{tab:progress-bounds}
\input{\ProgressReportDir/bounds.tex}
\end{table}

On the open subset, SAT has a better upper bound in
\ProgressOpenSATUB{} pairs, ILP in \ProgressOpenILPUB{}, and
\ProgressOpenEqualUB{} are tied. ILP has a stronger lower bound in
\ProgressOpenILPLB{} pairs; the remaining \ProgressOpenEqualLB{} are tied.
Maximum absolute differences reach \ProgressMaxUB{} label units for upper
bounds and \ProgressMaxLB{} for lower bounds. These discrepancies do not
support a claim that all differences are small or that the methods are
interchangeable. Gurobi proves more optima on this cohort, whereas SAT
provides better feasible spans on part of the unresolved subset.

There is \ProgressReturnedFeasible{} Gurobi observation recorded as
\FEAS{} after a normal worker return, rather than an external timeout.
Its native termination code was not saved, so its stopping reason cannot
be reconstructed. It is retained as \FEAS{} and is not counted as OPT.
The full observation tables preserve this distinction. These results
motivate reporting coverage and primal/dual bound quality together,
without inferring a general runtime or quality ranking over all graphs.
\FloatBarrier
